%titolo{Come al paguro Bernardo}
%autore{}
%album{}
%tonalita{DO}
%famiglia{Scout}
%gruppo{}
%momenti{}
%identificatore{come_al_paguro_bernardo}
%data_revisione{2026_07_1}
%trascrittore{Modanese Tobia}

\beginsong{Come al paguro Bernardo}

\beginverse
\[C]Quante volte è capi\[F]tato 
Di avere \[G7]avuto un aiuto inaspet\[C]tato
E dopo un po' di tempo aver ca\[F]pito 
Che tutto \[G7]questo non è stato mica un \[C]caso 
\[C]Quante volte è capi\[F]tato 
Di ve\[G7]der la gente in difficol\[C]tà 
Per ver\[C7]gogna o per orgoglio esage\[F]rato 
Una \[G7]mano agli altri mai la chiederà 
\endverse

\beginchorus
\[F]Come successe al pa\[C]guro Bernardo 
\[G7]Di professione ere\[A-]mita 
\[F]Gli fu salvata la \[C]vita 
Quando qual\[G7]cuno la conchiglia gli \[C]girò 
\endchorus
\chordsoff
\beginverse
Se a te è capitato
Di trovarti nel momento e al posto giusto
Come un efficiente angelo custode
Fare una buona azione eccezionale
Tirar via dai guai un malcapitato
Senza che s'accorga che... tu sei stato
E se ci riesci a farlo in gran segreto
Non c'è dono più grande che puoi far
\endverse


\beginverse
Se vuoi sapere chi è stato
A salvar la vita al povero eremita
Io di certo a nessuno lo dirò
Per paura che gli venga rivelato
Non si fanno mai le Buone Azioni
Solamente agli amici ed ai parenti
Per vantarsi e guadagnar apprezzamenti
Ed attendersi una qualche ricompensa
\endverse
\endsong
